Une mince nappe d'huile d'indice de réfraction \num{\v{n}} flotte sur une flaque d'eau (indice \num{1.33}). Son épaisseur est de \qty{\v{ep}}{nm}. On l'éclaire en lumière blanche et on la regarde de haut, presque à la verticale.

\begin{enumerate}[label=\alph*)]
    \item Quelle longueur d'onde du visible (\qtyrange{400}{700}{nm}) est la plus fortement réfléchie ? Justifiez en considérant les autres ordres.
    \item Quelle est la plus grande longueur d'onde complètement absente de la lumière réfléchie ?
    \item Quelle épaisseur minimale (non nulle) de la même huile donnerait la même couleur réfléchie qu'en a) ?
\end{enumerate}
